\documentclass[12pt,a4paper]{article}

\usepackage[czech]{babel}
\usepackage[T1]{fontenc}
\usepackage[utf8]{inputenc}

\usepackage{geometry}
\usepackage[most]{tcolorbox}

\geometry{margin=2.5cm}

\title{\textbf{Rozbor díla „R.U.R.“}}
\author{PAVPE}
\date{}

\begin{document}
\maketitle

%------------------------------------------------
\section*{Autor}

\begin{tcolorbox}[breakable]

\textbf{Autor:}
\begin{itemize}
\item Karel Čapek
\end{itemize}

\textbf{Název díla:}
\begin{itemize}
\item R.U.R. (Rossumovi univerzální roboti)
\end{itemize}

\textbf{Literární směr:}
\begin{itemize}
\item Česká meziválečná literatura (1918–1938), demokratický proud
\end{itemize}

\textbf{Století tvorby:}
\begin{itemize}
\item 20. století, 1. polovina
\end{itemize}

\textbf{Další autoři stejného směru + dílo:}
\begin{itemize}
\item Josef Čapek – Povídání o pejskovi a kočičce (1929)
\item Karel Poláček – Bylo nás pět (1946)
\item Eduard Bass – Cirkus Humberto (1941)
\end{itemize}

\textbf{Další autorova díla:}
\begin{itemize}
\item Boží muka (1917)
\item Továrna na absolutno (1922)
\item Válka s Mloky (1936)
\end{itemize}

\textbf{Specifický styl autora:}
\begin{itemize}
\item filozofická a společenská témata, varování před zneužitím techniky
\item využití dialogů, jednoduchý jazyk s hlubokou myšlenkou
\end{itemize}

\end{tcolorbox}

%------------------------------------------------
\section*{Charakteristika uměleckého textu}

\begin{tcolorbox}[breakable]

\textbf{Literární druh, žánr, forma:}
\begin{itemize}
	\item drama
	\item tragédie (s prvky sci-fi / antiutopie)
	\item dialogická forma (bez vypravěče)
\end{itemize}

\textbf{Typ vypravěče:}
\begin{itemize}
\item není přítomen (drama je tvořeno dialogy)
\end{itemize}

\textbf{Dominantní slohový postup:}
\begin{itemize}
\item dialogický
\end{itemize}

\textbf{Vysvětlení názvu díla:}
\begin{itemize}
\item R.U.R. = Rossumovi univerzální roboti
\item slovo „robot“ pochází z českého „robota“ (nucená práce)
\end{itemize}

\textbf{Aktuálnost díla:}
\begin{itemize}
\item dílo je aktuální i dnes (2026) – upozorňuje na nebezpečí technického pokroku a umělé inteligence
\end{itemize}

\textbf{Místo a čas děje:}
\begin{itemize}
\item Rossumův ostrov – továrna R.U.R.
\item neurčená budoucnost (hra napsána 1920, premiéra 1921)
\end{itemize}

\section*{Děj}

\subsection*{Prolog}
Děj začíná v kanceláři ředitele Dominova závodu R.U.R. Přichází Helena Gloryová, která chce roboty osvobodit, protože věří, že mohou mít duši.

Domin jí vysvětluje vznik robotů. Starý Rossum chtěl vytvořit umělého člověka, jeho synovec mladý Rossum výrobu zjednodušil a začal roboty vyrábět jako levnou pracovní sílu.

Helena se setkává s robotkou Sullou, kterou považuje za člověka. Postupně přicházejí další vedoucí pracovníci. Helena obhajuje lidskost, ostatní efektivitu.

Domin Helenu požádá o ruku.

\subsection*{1. dějství}
Helena žije s Dominem. Uvědomuje si důsledky existence robotů.

Z novin se dozvídá o poklesu porodnosti a využívání robotů ve válkách. Svět ztrácí přirozený řád.

Robot Radius odmítá pracovat a chce vládnout lidem. Helena spolupracuje s Dr. Gallem na „polidštění“ robotů.

\subsection*{2. dějství}
Roboti po celém světě povstávají. Lidé jsou v továrně obklíčeni.

Jedinou nadějí je tajemství výroby robotů. Helena však recept spálila.

Busman se snaží roboty uplatit, ale je zabit. Všichni lidé kromě Alquista umírají.

\subsection*{3. dějství}
Alquist zůstává posledním člověkem. Roboti ho nutí znovu objevit výrobu robotů.

Nedokáže je pitvat, protože není vědec. Objevuje však, že roboti Primus a Helena mají city.

Vidí v nich nového Adama a Evu – naději pro budoucnost.

\subsection*{Závěr}
Dílo varuje před ztrátou lidskosti a slepým technickým pokrokem.

\textbf{Smysl díla:}
\begin{itemize}
\item varování před zneužitím techniky a ztrátou lidskosti
\end{itemize}

\textbf{Pravděpodobný adresát:}
\begin{itemize}
\item společnost 20. století i současný člověk
\end{itemize}

\textbf{Kompozice:}
\begin{itemize}
\item chronologická (prolog + 3 dějství)
\end{itemize}

\textbf{Zařazení do kontextu autorovy tvorby:}
\begin{itemize}
\item jedno z nejvýznamnějších děl Karla Čapka, reakce na technický rozvoj po 1. světové válce
\end{itemize}

\textbf{Film nebo dramatizace:}
\begin{itemize}
\item dílo je často uváděno v divadlech
\item první televizní adaptace vznikly ve 20. století
\end{itemize}

\end{tcolorbox}

%------------------------------------------------
\section*{Úryvek k rozboru}

Helena: „Já věřím, že roboti jsou jako lidé. Že i oni mohou jednou cítit a myslet.“\\
\\
Domin: „To je omyl. Roboti nemají duši, jsou jen stroje vytvořené k práci.“\\
\\
Helena: „Ale proč by nemohli být víc než jen nástroje? Proč by nemohli být svobodní?“\\
\\
Domin: „Protože svět je potřebuje. Roboti pracují za lidi, zbavují je námahy a utrpení.“\\
\\
Helena: „A co když tím lidé ztrácejí sami sebe? Co když přestanou být lidmi?“\\
\\
Domin: „Naopak, lidstvo bude konečně šťastné. Nebude muset pracovat.“\\
\\
Helena: „Já nechci svět bez lidské práce. Práce dává životu smysl.“\\
\\
Domin: „To je starý názor. Nový svět bude lepší, efektivnější.“\\
\\
Helena: „Vy jste vytvořili něco, co nemůžete ovládat.“\\
\\
Domin: „Roboti jsou pod kontrolou.“\\
\\
Helena: „A co když se jednoho dne rozhodnou neposlouchat?“\\
\\
Domin: „To se nestane.“\\
\\
Helena: „Já tomu nevěřím.“\\

\begin{tcolorbox}[breakable]

\textbf{Atmosféra úryvku:}
\begin{itemize}
\item napjatá, filozofická
\item střet humanismu a technicismu
\end{itemize}

\textbf{Postavy v úryvku:}
\begin{itemize}
\item Helena Gloryová
\item Domin
\end{itemize}

\textbf{Charakteristika postav:}

\textbf{Hlavní postava}
\begin{itemize}
\item Helena – citlivá, humanistická, věří v lidskost robotů
\item zastává morální hodnoty
\end{itemize}

\textbf{Vedlejší postava}
\begin{itemize}
\item Domin – racionální, pragmatický, orientovaný na výkon a efektivitu
\end{itemize}

\textbf{Další postavy díla:}
\begin{itemize}
\item Dr. Gall
\item Alquist
\item Busman
\end{itemize}

\textbf{Vztahy mezi postavami:}
\begin{itemize}
\item názorový konflikt, později partnerský vztah
\end{itemize}

\textbf{Zařazení úryvku do děje:}
\begin{itemize}
\item prolog – seznámení s hlavní myšlenkou díla
\end{itemize}

\textbf{Jazykové prostředky:}
\begin{itemize}
\item metafora – „lidé ztrácejí sami sebe“
\item opakování – důraz na otázky Heleny
\item kontrast – humanita vs technika
\end{itemize}

\end{tcolorbox}

%------------------------------------------------
\section*{Názor na dílo}
Dílo R.U.R. na mě působilo silně především svou myšlenkou. Nejvíce mě zaujal závěr, který odkazuje na biblický motiv Adama a Evy. Text je místy náročnější na sledování kvůli většímu množství postav, ale hlavní myšlenka je velmi aktuální i dnes. Dílo ukazuje nebezpečí technického pokroku bez morálních hodnot.

	\end{document}
